\باب{کوشی کا  مسئلہ اوسط قیمت اور قاعدہ لھوپیٹال  ے کا پائیدار روپ}\شناخت{ضمیمہ_ہ}
اس ضمیمہ میں قاعدہ    لھوپیٹال (حصہ \حوالہ{حصہ_ماورائی_قاعدہ_لھوپیٹال}، مسئلہ \حوالہ{مسئلہ_ماورائی_لھوپیٹال_دوم}) کے پائیدار  روپ کی محدود حد صورت کا ثبوت پیش کیا جاتا ہے۔

\حصہء{قاعدہ    لھوپیٹال(پائیدار روپ)}
فرض کریں     کھلے وقفہ \عددی{(a,b)} پر، جس میں نقطہ \عددی{x_0} پایا جاتا ہے ، تفاعلات \عددی{f} اور \عددی{g}  دونوں قابل تفرق ہیں اور درج ذیل ہے۔
\[f(x_0)=g(x_0)=0\]
مزید فرض کریں   پورے \عددی{(a,b)} میں  ،ممکنہ طور پر   ماسوائے  نقطہ \عددی{x_0} کے، ہر نقطے پر \عددی{g'\ne 0}  ہے۔تب ذیل ہو گا، بشرطیکہ  دائیں ہاتھ  پر حد موجود ہو۔
\begin{align}\label{مساوات_ضمیمہ_پائیدار_روپ}
\lim_{x\to x_0}\frac{f(x)}{g(x)}=\lim_{x\to x_0}\frac{f'(x)}{g'(x)}
\end{align}

قاعدہ    لھوپیٹال کے پائیدار روپ کا ثبوت کوشی کے مسئلہ اوسط قیمت  پر منحصر ہے، جو ایسا اوسط قیمت مسئلہ ہے جو  ایک کے بجائے دو تفاعلات کے متعلق ہے۔ ہم   مسئلہ کوشی   ثابت کرنے کے بعد دکھاتے ہیں  اس سے قاعدہ    لھوپیٹال  کیسے اخذ ہوتا ہے۔

\ابتدا{مسئلہءء}\موٹا{کوشی کا اوسط قیمت مسئلہ}\\
فرض کریں \عددی{[a,b]} پر تفاعلات \عددی{f} اور \عددی{g} استمراری اور پورے  \عددی{(a,b)} پر  قابل تفرق ہیں، اور پورے \عددی{(a,b)} پر \عددی{g'\ne 0} ہے۔تب \عددی{(a,b)} میں  ایسا عدد \عددی{c} موجود ہو گا جس پر درج ذیل مطمئن ہوتا ہو۔
\begin{align}\label{مساوات_ضمیمہ_کوشی_مسئلہ}
\frac{f'(c)}{g'(c)}=\frac{f(b)-f(a)}{g(b)-g(a)}
\end{align}
  اس میں \عددی{g(x)=x}  رکھنے سے  مسئلہ اوسط قیمت  کا سادہ روپ (حصہ \حوالہ{حصہ_استعمال_مسئلہ_اوسط_قیمت}، مسئلہ \حوالہ{مسئلہ_استعمال_اوسط_قیمت}) حاصل ہو گا۔
\انتہا{مسئلہءء}

\ابتدا{ثبوتءء}\موٹا{کوشی کے  اوسط  قیمت  مسئلے کا ثبوت}\\
ہم حصہ \حوالہ{حصہ_استعمال_مسئلہ_اوسط_قیمت} کے مسئلہ اوسط قیمت کا اطلاق دو مرتبہ کرتے ہیں۔اس کو استعمال کر کے ہم پہلے دکھاتے ہیں کہ \عددی{g(a){\ne} g(b)} ہو گا۔اگر \عددی{g(b)}  اور \عددی{g(a)} برابر ہوں تب \عددی{a} اور \عددی{b} کے بیچ کسی \عددی{c} پر  مسئلہ اوسط قیمت درج ذیل دیتا
\[g'(c)=\frac{g(b)-g(a)}{b-a}=0\]
تاہم \عددی{(a,b)} میں \عددی{g'(x)\ne 0} ہے لہٰذا ایسا ممکن نہیں۔

اس کے بعد مسئلہ اوسط قیمت کا اطلاق  ہم ذیل تفاعل پر  کرتے ہیں۔
\[F(x)=f(x)-f(a)-\frac{f(b)-f(a)}{g(b)-g(a)}[g(x)-g(a)]\]
یہ تفاعل وہاں استمراری اور قابل تفرق ہے جہاں \عددی{f} اور \عددی{g} استمراری اور قابل تفرق ہیں، اور \عددی{F(b)=F(a)=0} ہے۔لہٰذا \عددی{a} اور \عددی{b} کے بیچ ایسا   عدد \عددی{c} موجود ہو گا جس پر \عددی{F'(c)=0} ہو گا۔تفاعلات \عددی{f} اور \عددی{g} کی صورت میں یہ حقیقت درج ذیل لکھی جا سکتی ہے
\[F'(c)=f'(c)-\frac{f(b)-f(a)}{g(b)-g(a)}[g'(c)]=0\]
جس سے ذیل لکھا جا سکتا ہے جو مساوات \حوالہ{مساوات_ضمیمہ_کوشی_مسئلہ} ہے۔
\[\frac{f'(c)}{g'(c)}=\frac{f(b)-f(a)}{g(b)-g(a)}\]
\انتہا{ثبوتءء}

\ابتدا{ثبوتءء}\موٹا{قاعدہ    لھوپیٹال کے پائیدار روپ کا ثبوت }\\
ہم اول  \عددی{x\to x_0^+}  کی صورت میں مساوات \حوالہ{مساوات_ضمیمہ_پائیدار_روپ}  کی تصدیق کرتے ہیں۔یہی طریقہ \عددی{x\to x_0^-} کے لئے بھی  استعمال  ہو گا، اور دونوں کو   ملا کر نتیجہ    حاصل ہو گا۔

فرض کریں \عددی{x_0} کے دائیں جانب \عددی{x} پایا جاتا ہے۔یوں \عددی{g'(x)\ne 0} ہو گا لہٰذا \عددی{x_0} سے لیکر \عددی{x}  تک کے بند وقفے پر کوشی کے  مسئلہ اوسط قیمت کا اطلاق ہو گا۔اس سے \عددی{x_0} اور \عددی{x} کے بیچ ایسا عدد \عددی{c} حاصل ہو گا جس پر درج ذیل مطمئن ہو گا،
\[\frac{f'(c)}{g'(c)}=\frac{f(x)-f(x_0)}{g(x)-g(x_0)}\]
تاہم \عددی{f(x_0)=g(x_0)=0}  ہے لہٰذا درج ذیل   ہو گا۔
\[\frac{f'(c)}{g'(c)}=\frac{f(x)}{g(x)}\]
جیسے جیسے \عددی{x} کی قیمت \عددی{x_0} کی قیمت کے قریب پہنچتی ہے، چونکہ \عددی{c} خود \عددی{x} اور \عددی{x_0} کے بیچ پایا جاتا ہے لہٰذا   \عددی{c}  کی قیمت \عددی{x_0} کو پہنچے گی۔یوں درج ذیل ہو گا۔
\[\lim_{x\to x_0^+}\frac{f(x)}{g(x)}=\lim_{c\to x_0^+}\frac{f'(c)}{g'(c)}=\lim_{x\to x_0^+}\frac{f'(x)}{g'(x)}\]

یوں   قاعدہ    لھوپیٹال   کا ثبوت اس صورت  کے لئے مکمل ہوا جب \عددی{x} کی قیمت \عددی{x_0} کی قیمت  تک اوپر سے پہنچتی ہو۔ جہاں \عددی{x} کی قیمت \عددی{x_0} کی قیمت تک نیچے  سے  پہنچتی ہو وہاں ہم \عددی{x<x_0}   لے کر  کوشی کے اوسط قیمت مسئلہ کا اطلاق وقفہ \عددی{[x,x_0]}پر کرتے ہوئے ثبوت حاصل کرتے ہیں۔
\انتہا{ثبوتءء}

